% !TeX root = ../main.tex
% !TeX encoding = utf8

\chapter{Diseño de puestos y superestructura}
% Responsable: David Bacas Posadas
% Secciones del índice:
%   4.1 Diseño de puestos (preparación y adoctrinamiento — continuación)
%   4.2 Diseño de la superestructura

\section{Diseño de puestos: preparación}

% 4.1 — Preparación
%
%   - ¿Qué titulación o formación previa exige cada puesto?
%   - ¿La formación requerida otorga autonomía al trabajador o está igualmente regulado?
%   - Identificar los puestos que requieren preparación profesional (universitaria o técnica).
%   - Describir la formación interna que ofrece CaixaBank (tipo, frecuencia, obligatoriedad).

La preparación es el parámetro de diseño mediante el cual la organización asegura que el empleado posee los conocimientos y habilidades necesarios antes de ocupar el puesto. En la sucursal de Motril, el nivel de preparación exigido varía sustancialmente según el perfil del puesto, aunque se enmarca en las exigencias globales de CaixaBank.

Para los perfiles de asesoramiento experto (Gestor Premier, Gestor de Banca Personal y Gestor de Negocio), la entidad exige una alta cualificación previa, generalmente una titulación universitaria superior en áreas como Economía o Administración y Dirección de Empresas (ADE). Por el contrario, para el puesto de Gestor Comercial —más orientado a tareas operativas y transaccionales— los requisitos académicos iniciales pueden ser más flexibles, aceptándose perfiles de Formación Profesional (Grado Superior en Administración y Finanzas).

Una vez incorporados a la sucursal, todos los empleados acceden a ``Virtaula'', el robusto campus virtual de formación continua de la entidad. Esta formación interna se estructura en tres grandes bloques:
\begin{itemize}
    \item \textbf{Formación Normativa Obligatoria:} Es crucial para los gestores especializados de Motril. Destacan las certificaciones exigidas por el regulador, como la certificación MiFID II (Inversiones, 150 horas) para los gestores Premier y de Banca Personal, o la certificación LCCI para crédito inmobiliario. La superación de estos exámenes es condición \textit{sine qua non} para ejercer el asesoramiento.
    \item \textbf{Formación para la Retribución Variable:} Cursos obligatorios para toda la plantilla sobre prevención de blanqueo de capitales, código ético y ciberseguridad, indispensables para el cobro de incentivos.
    \item \textbf{Formación de Desarrollo y Liderazgo:} Programas de desarrollo orientados a los puestos directivos de la sucursal (Directora y Subdirectora), enfocados en potenciar sus habilidades de gestión de equipos y visión estratégica.
\end{itemize}

Si bien esta alta preparación dota a los gestores especializados de las capacidades para desempeñar su función, la autonomía en la sucursal está fuertemente acotada por los protocolos informáticos y regulatorios. No obstante, el modelo corporativo (modelo AHEAD) promueve cierto grado de empoderamiento, otorgando a los gestores mayor margen de maniobra en la gestión emocional y relacional con los clientes de Motril.

\section{Diseño de puestos: adoctrinamiento}

% 4.1 — Adoctrinamiento
%
%   - ¿Existe una cultura corporativa o valores fuertemente integrados en CaixaBank?
%   - ¿Cómo se inculcan esos valores (acogida, cursos, ejemplo directivo...)?
%   - ¿Se valora la afinidad cultural en el proceso de selección?

El adoctrinamiento es el proceso mediante el cual la organización socializa a sus miembros para que adopten sus valores y creencias. En CaixaBank, y por extensión en la sucursal de Motril, la cultura corporativa actúa como un fuerte mecanismo de cohesión que guía el comportamiento, supliendo el margen de maniobra que dejan los protocolos. Esta cultura se basa en pilares fundamentales como la vocación de servicio, la colaboración y la agilidad.

Estos valores se inculcan desde el proceso de selección corporativo, donde se evalúa la afinidad del candidato con el ``ADN'' de la entidad mediante entrevistas conductuales. Una vez asignados a la sucursal, el adoctrinamiento continúa mediante programas de acogida (\textit{onboarding}). Los nuevos empleados cuentan con tutores o \textit{buddies} dentro de la propia oficina de Motril para facilitar su integración. 

Asimismo, el adoctrinamiento se ejerce a nivel local de forma descendente y ejemplarizante: se espera que la Directora y la Subdirectora actúen como referentes culturales de estos valores, fomentando un clima de colaboración en el equipo que mitigue la rigidez del trabajo diario. Además, se incentiva la participación de los empleados de la sucursal en acciones de voluntariado a través de la Fundación ``la Caixa'', fortaleciendo así su vínculo y compromiso social con el entorno local de Motril.

\section{Diseño de la superestructura}

% 4.2 — Superestructura
%
% Para cada unidad o departamento identificado:
%
%   Tamaño:
%     - Número de personas que integran cada unidad.
%     - Ámbito de control de cada directivo (cuántas personas supervisa directamente).
%
%   Base de agrupación:
%     - ¿Agrupación funcional (por tipo de tarea) o por mercado (por tipo de cliente)?
%     - Justificar con ejemplos concretos de la estructura de la sucursal.

La superestructura se refiere a cómo se agrupan los puestos en unidades y cuál es el tamaño de estas. El diseño de la sucursal de Motril responde a un modelo de especialización comercial que busca maximizar la eficiencia en la atención presencial.

\subsection*{Tamaño de la unidad y ámbito de control}
La sucursal de Motril analizada constituye una unidad de tamaño reducido. Aunque la plantilla exacta está pendiente de confirmación, la estructura base se articula en torno a \textbf{seis roles o perfiles principales} (dos puestos directivos y cuatro perfiles operativos). Este tamaño de unidad permite una gestión directa y un clima de gran proximidad entre sus miembros.

El ámbito de control (\textit{span of control}) recae principalmente sobre la Directora de la sucursal, quien supervisa directamente a todo el equipo. Dado el tamaño reducido de la oficina, la Directora ejerce un tramo de control estrecho sobre la Subdirectora y el equipo de gestores operativos. Este número limitado de perfiles a su cargo permite una supervisión directa efectiva y un flujo de comunicación continuo, facilitando que la Directora se involucre no solo en el control administrativo, sino también en el apoyo comercial y la dinamización del equipo.

\subsection*{Base de agrupación}
Históricamente, las oficinas bancarias utilizaban una base de agrupación funcional (por ejemplo, separando al personal de caja del personal comercial). Sin embargo, en la sucursal de Motril se evidencia claramente una \textbf{base de agrupación por mercado}, concretamente segmentada por el \textbf{tipo de cliente}. 

Esta agrupación permite que cada perfil conozca en profundidad las necesidades específicas del segmento que atiende. La estructura de la oficina materializa esta agrupación en sus cuatro puestos operativos:
\begin{itemize}
    \item \textbf{Gestor/a de Banca Personal:} Agrupa a los clientes particulares con necesidades de planificación patrimonial de nivel medio-alto.
    \item \textbf{Gestor/a Premier:} Agrupa a los clientes con un perfil de inversión más sofisticado y mayor patrimonio.
    \item \textbf{Gestor/a de Negocio:} Centraliza la atención a pequeñas empresas, comercios y autónomos de Motril, gestionando sus particularidades de crédito y tesorería.
    \item \textbf{Gestor/a Comercial:} Agrupa a los clientes particulares de segmento estándar para las gestiones cotidianas y transaccionales.
\end{itemize}

En conclusión, el diseño organizativo de la sucursal de Motril sacrifica la polivalencia total de su plantilla en favor de una alta especialización orientada al mercado. Esta base de agrupación por cliente, sostenida por la centralización tecnológica que asume las tareas puramente administrativas, garantiza que la sucursal se enfoque casi en exclusiva en el asesoramiento y la actividad comercial de valor añadido.

\endinput
% -------------------------------------------------------------------
% FIN DEL CAPÍTULO — David Bacas Posadas
% -------------------------------------------------------------------
